% ====================================================================
% §6 평형 peak — |I|→0 기준선 (N6)
% ====================================================================
\section{평형 peak --- $|I|\to0$ 기준선}\label{sec:eqpeak}
평형 진행률에서 곧바로 얻어지는 것이 평형 peak 모양 $\xi_\eq(1-\xi_\eq)/w$ 이다. 이것이 전류가 없을 때($|I|\to0$)의
기준선이며, 동역학 꼬리가 무시될 만큼 작을 때 곡선이 이 종으로 직접 환원되는 항이다.

\textbf{(a) 출발 --- 전하 보존식.} 관측 $\dd Q/\dd V$ 는 전하 보존식 $Q_\cell q=Q_\bg(V_n)+\sum_jQ_j\xi_j$ 를 궤적
$q$ 로 미분해 닫힌다 --- 양변을 $q$ 로 미분하면 $Q_\cell=C_\bg\,\dd V_n/\dd q+\sum_jQ_j\,\dd\xi_j/\dd q$ 이고,
$\dd V_n/\dd q$ 로 나누면 $\dd Q/\dd V_n=C_\bg+\sum_jQ_j\,\dd\xi_j/\dd V_n$ 이다. \textbf{(b) 연산 --- logistic
미분의 종 항등식.} 평형 추종이면 $\dd\xi_j/\dd V_n$ 에 평형 기울기가 들어가고, $z\equiv\sigma_d(V-U_j^{\,d})/w_j$ 로
$\xi_\eq=(1+e^{-z})^{-1}$ 를 미분하면
\begin{equation}
\frac{\dd\xi_\eq}{\dd z}=\frac{e^{-z}}{(1+e^{-z})^2}
=\underbrace{\frac{1}{1+e^{-z}}}_{\xi_\eq}\cdot\underbrace{\frac{e^{-z}}{1+e^{-z}}}_{1-\xi_\eq}
=\xi_\eq(1-\xi_\eq)
\label{eq:belliden}
\end{equation}
--- $\xi(1-\xi)$ 는 $\xi=\tfrac12$ 에서 최대 $\tfrac14$ 인 종이다. \textbf{(c) 중간식 --- 연쇄율.} 연쇄율
$\dd z/\dd V=\sigma_d/w_j$ 를 곱하면 $\dd\xi_{\eq,j}/\dd V=\sigma_d\,\xi_\eq(1-\xi_\eq)/w_j$. \textbf{(d) 박스.}
이를 보존식 미분에 넣으면 전이 $j$ 의 평형 peak 이 닫힌다:
\begin{equation}
\frac{\dd\xi_{\eq,j}}{\dd V}=\frac{\sigma_d\,\xi_{\eq,j}(1-\xi_{\eq,j})}{w_j}
\;\Longrightarrow\;
\boxed{\;\Big(\frac{\dd Q}{\dd V}\Big)^\eq_j
=Q_j\,\frac{\xi_{\eq,j}(1-\xi_{\eq,j})}{w_j}
=Q_j\,\Big|\frac{\dd\xi_{\eq,j}}{\dd V}\Big|\;.}
\label{eq:eqpeak}
\end{equation}
이 값이 전이별로 합산되어 배경과 함께 $|I|\to0$ 극한을 이룬다 --- 단 $\gamma_j\ne0$ 이면 이 극한은 분기
중심 $U_j^{\,d}$ 에 남는 히스테리시스 잔존 극한이고, 분기 없는 가역 기준선($U_j$ 중심)과는
$\gamma_j\to0$ 에서만 일치한다(히스테리시스는 열역학적 --- \S\ref{sec:hys}). 부호 --- $\sigma_d$ 가 미분에서 한 번 들어오지만 peak 모양은
절댓값 $|\dd\xi/\dd V|=\xi_\eq(1-\xi_\eq)/w$ 로 양수이고 방향 불변이다($\xi(1-\xi)\ge0$). 세 양 --- 위치 $=$ 중심
$U_j^{\,d}$, 순높이 $=Q_j/(4w_j)$($\xi=\tfrac12$ 에서 $\xi(1-\xi)=\tfrac14$), 배경 차감 면적 $=Q_j$
($\int_0^1\dd\xi=1$, 식~\eqref{eq:belliden} 의 종이 $\dd\xi_\eq$ 의 치환적분이라 면적 $1$) --- 가 이 한 식에서 읽힌다.

\textbf{선택 확장 --- skew-logistic 비대칭 종(opt-in).} 식~\eqref{eq:eqpeak} 의 종은
$\xi_\eq\!\leftrightarrow\!1-\xi_\eq$ 교환에 대해 대칭이라 그 자체로는 좌우 비대칭을 담지 못한다.
진행률 좌표를 식~\eqref{eq:xieq} 의 logistic 을 밑으로 하는 거듭제곱
$\xi^{(\alpha)}_{\eq,j}\equiv\big[\xi_{\eq,j}\big]^{\alpha_j}$ ($\alpha_j>0$)로 재모수화하면 --- 이 좌표도
$V$ 에 대해 여전히 단조 $0\!\to\!1$ 이다 --- 위 (b)(c) 와 똑같은 연쇄율 계산이 비대칭 종을 준다:
\begin{equation}
\boxed{\;\Big(\frac{\dd Q}{\dd V}\Big)^\eq_j
=Q_j\,\frac{\alpha_j}{w_j}\,\big[\xi_{\eq,j}\big]^{\alpha_j}\big(1-\xi_{\eq,j}\big)\;,
\qquad \xi^{(\alpha)}_{\eq,j}=\big[\xi_{\eq,j}\big]^{\alpha_j}\;.}
\label{eq:skewpeak}
\end{equation}
$\alpha_j=1$ 이면 $[\xi_\eq]^1(1-\xi_\eq)/w_j$ 로 식~\eqref{eq:eqpeak} 를 \emph{정확히} 회수한다(구현도
이 경로가 부동소수점까지 동일 --- 부록 B). 위 \emph{세 양}은 이렇게 일반화된다. \emph{면적}은
$\int_0^1\dd\xi^{(\alpha)}_\eq=1$ 이라 $\alpha_j$ 에 무관하게 $Q_j$ 로 \emph{보존}된다. 그러나
\emph{정점}은 $\xi_{\eq,j}=\alpha_j/(\alpha_j+1)$ 에서 걸리므로, 그 \emph{위치}와 \emph{높이}가 함께 옮겨간다:
\begin{equation}
V^\ast_j-U_j^{\,d}=\sigma_d\,w_j\ln\alpha_j\;,\qquad
\Big(\frac{\dd Q}{\dd V}\Big)^\eq_{j}\Big|_{V^\ast_j}
=\frac{Q_j}{w_j}\Big(\frac{\alpha_j}{\alpha_j+1}\Big)^{\alpha_j+1}
\label{eq:skewapex}
\end{equation}
($\alpha_j=1$ 에서 각각 $0$ 과 $Q_j/(4w_j)$ --- 위 세 양과 일치; 수치 확인). 곧 $\alpha_j\ne1$ 에서는
\emph{정점이 더는 전이 중심이 아니다} --- 중심 $U_j^{\,d}$ 의
\emph{정의}(식~\eqref{eq:xieq} 의 logistic 중심, $\xi_\eq=\tfrac12$)는 불변이지만 정점에서 직접 읽어낼 수
없다($w_j\!\approx\!25.7$ mV 에서 $\alpha_j=2$ 면 $+17.8$ mV 이동). 폭도 함께 변한다 --- 반높이 폭은
$\alpha_j=0.5/1/2/4$ 에서 각각 $4.45/3.53/3.02/2.74$ 배 $w_j$ 로 \emph{$\alpha_j$ 와 함께 단조 감소}하고
($\alpha_j<1$ 은 기준보다 오히려 넓다), 정점 좌우 반폭비는 $0.80/1.00/1.18/1.30$ 이 된다(수치 확인).
그러므로 $\alpha_j$ 는 순수한 비대칭 손잡이가 아니라
\emph{위치$\cdot$폭$\cdot$비대칭이 한데 묶인} 손잡이이며, 그 때문에 $w_j$ 및
$L_{V,j}$(\S\ref{sec:lag})와 식별 축퇴를 일으킨다 --- 가드는 \S\ref{sec:inputs} 가 세운다.
종은 $\alpha_j>0$ 전역에서 $C^\infty$ 이고, $\xi^{(\alpha)}_\eq$ 가 단조라
\S\ref{sec:lag}--\S\ref{sec:tail} 의 인과 꼬리와 전하 보존 반전(식~\eqref{eq:implicit})이 그대로 성립한다.
\emph{지위} --- $\alpha_j$ 는 전류$\cdot$과전압과 무관한 \emph{평형} 형상 파라미터(현상학적$\cdot$tier C)로
새 상(phase)이 아니며, 물리적 동기는 order--disorder 엔트로피 스텝$\cdot$조성 의존 $\Omega(x)$ 류의 조성
비대칭이다. 기본값은 $\alpha_j=1$(전이 입력에 없으면 현행과 동일 --- opt-in)이고, 바로 아래 배경 상자의
감수율 \emph{유도 서식}($\mathrm{var}_j=M_j\theta_j(1-\theta_j)$ 의 독립-자리 분산)은 $\alpha_j=1$ 기준이다.

\begin{bgbox}
\textbf{평형 peak 은 등온 감수율 --- fluctuation--dissipation 읽기.} 식~\eqref{eq:eqpeak} 의 평형 peak 은 곡선맞춤 종이
아니라 전극에 저장된 전하의 \emph{등온 입자수 감수율}(isothermal particle-number susceptibility) 그
자체다. Part 0 이 자리 하나에서 fluctuation--response 관계 $\mathrm{var}(n)=\theta(1-\theta)$(식~\eqref{eq:sm-flucres})를,
다클래스에서 $\partial\langle N\rangle/\partial\mu=\beta\,\mathrm{var}(N)$(식~\eqref{eq:sm-mc-fluc})를 세웠고,
$\langle N\rangle=\sum_jM_j\theta_j$(식~\eqref{eq:sm-mc-occ})였다. 이 fluctuation 을 관측 $\dd Q/\dd V$ 에 곧장
잇는다. 이상 폭($w_j=RT/F$, 곧 $n_j=1$)에서 식~\eqref{eq:eqpeak} 의 전이 몫은 $Q_j=(F/N_A)M_j$ 와
$\theta(1-\theta)=\xi(1-\xi)$ 로
\[
\Big(\frac{\dd Q}{\dd V}\Big)^\eq_j
=\frac{Q_j\,\xi_{\eq,j}(1-\xi_{\eq,j})}{w_j}
=\frac{F^2}{N_A RT}\,M_j\,\theta_j(1-\theta_j)
=\frac{F^2}{N_A RT}\,\mathrm{var}_j(N),
\]
전 전이를 더하고 $\mathrm{var}(N)=\kB T\,\partial\langle N\rangle/\partial\mu$(식~\eqref{eq:sm-mc-fluc})와
$\kB/R=1/N_A$ 를 쓰면
\[
\boxed{\;\Big(\frac{\dd Q}{\dd V}\Big)^\eq
=\frac{F^2}{N_A RT}\,\mathrm{var}(N)
=e^2\,\frac{\partial\langle N\rangle}{\partial\mu}\;}
\qquad(e\equiv F/N_A,\ \text{이상 폭}),
\]
곧 \emph{평형 $\dd Q/\dd V$ 는 저장 전하의 fluctuation 스펙트럼}이다. $\mu\leftrightarrow V$
결선(식~\eqref{eq:sm-eqcond})이 감수율 $\partial\langle N\rangle/\partial\mu$ 를 전위축 응답
$\partial Q/\partial V$ 로 옮긴 것이 이 항등식 --- fluctuation--dissipation 관계의 전기화학판이며, 자리당 $e^2$ 과 몰당
$F^2/N_ART$ 는 같은 환산의 두 표기다($e$ 는 자리당 전하, $F=N_Ae$). \emph{봉우리 하나가 fluctuation 하나}이고,
곡선 전체가 평형 fluctuation 의 지문이다 --- Part 0 이 fluctuation positivity $\partial\langle N\rangle/\partial\mu>0$ 으로
유일한 근을 보증한 바로 그 $\mathrm{var}(N)$(식~\eqref{eq:sm-mc-fluc})이, 여기서는 peak 높이로 관측된다.
\begin{verifybox}
검산 --- (i) \emph{합규칙(면적 $=$ 용량).} $\int(\dd Q/\dd V)^\eq_j\,\dd V=Q_j\!\int_0^1\dd\xi=Q_j$ 라
감수율의 전위 적분이 전이 총 전하다 --- fluctuation 스펙트럼의 합규칙이 곧 이동 가능한 전하량(식~\eqref{eq:eqpeak}
면적 읽기와 동일). (ii) \emph{정점.} $\xi=\tfrac12$ 에서 $\mathrm{var}$ 가 최대라 순높이 $Q_j/(4w_j)$ --- fluctuation 이
가장 큰 반충전에서 감수율이 최대다. (iii) \emph{폭 이중지위 가드.} 이 감수율 항등의 \emph{유도 서식}($\mathrm{var}_j=M_j\theta_j(1-\theta_j)$
의 독립-자리 분산과 $w_j=RT/F$)은 이상 격자기체($\Omega_j=0$)이자 \emph{대칭 종($\alpha_j{=}1$)}에서 정확하다
--- skew($\alpha_j\ne1$, 식~\eqref{eq:skewpeak})는 이 대칭 분산을 깨어 항등이 형상 파라미터로 격하되고, $n_j=1$ 이어도
$\Omega_j\neq0$ 이면(본문 기본 상태 $\Omega_j>2RT$ 포함) 이 서식은 근사가 되고, fluctuation--response 등식
$\mathrm{var}(N)=\kB T\,\partial\langle N\rangle/\partial\mu$ 자체만 대정준 일반으로 남는다. 두-상 전이의
현상학적 자유 폭($w_j$ 가 broadening 을 흡수 --- \S\ref{sec:width}$\cdot$\S\ref{sec:broadening})에서는
peak 이 그 감수율을 폭 $n_j$ 배로 펴고 높이 $1/n_j$ 배로 누른 상이라 문자 그대로의 감수율 \emph{크기}가
재척도된다(형상 근거는 유지되나 두-상 폭은 평형 예측이 아니다). (iv)
\emph{배경 제외.} 비전이 배경 $C_\bg$(식~\eqref{eq:sum})는 클래스 밖 저장이라 이 fluctuation 항에 들지 않는다
(그 $C_\bg$ 를 상수로 두는 것 자체가 \emph{피팅 창-국소} 근사임은 \S\ref{sec:sum} 이 명기한다 --- 실측 배경은
창을 넓히면 단조 감쇠한다).
\end{verifybox}
\end{bgbox}
% 자산(v1.0.22 SM2): [V22-SM2-A 평형 peak=등온 감수율 항등(dQ/dV=e²∂⟨N⟩/∂μ·합규칙·폭 이중지위 가드)]


이 식이 쓰이는 조건은 지연 길이가 peak 폭에 비해 무시될 만큼 작을 때다 --- \S\ref{sec:lag}(N7)이 그 지연
길이 $L_V$ 를 세우고, $L_{V,j}\ll w_j$ 면(또는 동역학 입력이 없으면) 식~\eqref{eq:eqpeak} 의 평형 종으로
매끈히 환원된다(식~\eqref{eq:tail-limit} 의 $L_V\to0$ 극한 --- 초기 상태 수치로는 $L_{V,j}\sim10^{-10}$--$10^{-8}$ V
대 $w_j$ 수십 mV, \S\ref{sec:sum}). 이것이 $|I|\to0$ 환원이다.
그 전에, 다음 절(\S\ref{sec:broadening})이 이 평형 종이 \emph{두-상 전이}에서 어떻게 실측 종으로 펴지는지 ---
곧 폭 $w_j$ 의 두-상 지위 --- 를 모델 안의 양들로 닫는다.

% 자산: [A-123] [A-124] [A-125] [A-126] [A-127] [A-128] [A-129]
% 자산(v1.0.25 C1): [V1025-C1-A skew-logistic 비대칭 종 eq:skewpeak(α=1→eq:eqpeak 정확 회수·면적 Q_j 보존)]
%                   [V1025-C1-B 정점 이동·높이 재척도 eq:skewapex(V*−U=σ_d·w·lnα · 높이=(Q/w)(α/(α+1))^{α+1})]
